\section{Solving the encoding problem}
\label{introduction:encodingSolution}

I have painted an admittedly bleak picture of the difficulties in studying
cell signaling. It should be clear from this discussion that experimental
design in cell signaling is quite difficult, and interpretation of experiments
must be done with extreme care. But can we remove, or at least reduce,
some of the difficulties discussed above?


\subsubsection{Obtaining meaningful $S$ and $R$}

Perhaps the most problematic of the issues discussed is that
of not knowing which signals $\vec{S}$ a cell is interested in 
nor which responses $\vec{R}$ encode those signals. One can begin to address
this by considering different variations of the inputs (e.g. changing
concentrations or treatment durations) and outputs
(e.g. phosphorylation states, live-cell measurements of the same
cell over time, fold-change in concentration, total change in concentration).
Further discussion of different output measurements, in the particular
context of single-cell fluorescence microscopy, can
be found in \ar{imaging:introduction}.


Assuming that one could identify a series of potential input signals
and response metrics, it is not obvious how to go about identifying
which are the more accurate estimates of $\vec{S}$ and $\vec{R}$.
An interesting
approach would be to perform measurements of information content
between putative combinations of $S$ and $R$, for example using the
mutual information metric \cite{Cheong2011}.
Under the assumption
that the cell is encoding signals in the most informative way
possible, the $R$ and $S$ choices that maximizes mutual information can then
be considered to be the best approximation
of the encoding that the cell uses. This would require either high
accuracy in measurement or precise knowledge of measurement error. 
In \ar{insulation:introduction} I test multiple reagents and readouts
to ensure that they have similar information content, and I discuss
this approach in more detail in
\ar{imaging:introduction}.


In an example of this approach, research groups measured the
information content of transcription factor gradients across
\fly\ embryos, with the question of whether enough information
was present to specify the location of all nuclei along the embryo.
While each transcription factor had low positional information content when taken alone,
in combination the factors did have enough information to specify each nucleus
position with high accuracy \cite{Gregor2007,Dubuis2011,Dubuis2013}.


As with the example of human speech above, 
it is important to remember that low information content of a
single $S$ does not
necessarily mean that it is an \i{incorrect}
encoding. It may alternatively signify an \i{incomplete} encoding,
and that that other unmeasured signal properties need to also be considered.
Importantly, incomplete encodings can be good enough for many
experimental biology questions.



\subsubsection{Compensating for cellular variability}


To address issues stemming from cellular variability, the
straightforward solution is to directly measure its effects on the
experimental relationship between the chosen $S$ and $R$. Measurements
of $R$ can be performed on a single-cell basis, for example by
microscopy (as in this dissertation) or by flow sorting.
The distributions of single-cell values can then be checked for properties,
such as multi-modality, that would suggest the presence of multiple
cellular states. (In my work (\ar{insulation:introduction}),
I verified that each measurement generated unimodal single-cell distributions.)


In the case that different cellular states do exists, each cell
can be grouped into statistically distinct subpopulations by measured
phenotype \cite{Loo2010,Singh2010}. Each subpopulation can then be
tested separately to see if each has the same $f(S)=R$ relationship.
For example, in \ar{insulation:introduction} I test how cell cycle phase
affects measurements of inter-pathway crosstalk. Further, if live-cell markers
are able differentiate subpopulations, then cells can be physically
sorted and experimented on separately.


Importantly, the absence of subpopulations
along the dimension of the measured response $R$ does not imply that
all cells are the same. Rather, it implies that we cannot claim that they
are different. Conversely, the presence of subpopulations in one
dimension does not imply existence of the same subpopulations with respect
to other dimensions (see \ar{fig:introduction:subpops}). In other
words, the presence or absence of cellular subpopulations as measured
by one readout is insufficient evidence to make a claim about whether $R$
is being distorted by the presence of subpopulations. Such a connection
must be explicitly tested.


  \begin{figure}[!bt]
  \centering
  \includegraphics[width=2.5in]{FIGS/introduction/subpops.pdf}
  {\singlespacing 
  \caption[Cellular subpopulations are phenotype-dependent]
          { The presence of subpopulations in one measurement dimension
            does not imply subpopulations in another dimension;
            subpopulations are a phenotype-dependent property.
            (\b{a}) Cells show bi-modality in their total DNA content
            (as measured by total Hoechst fluorescence), but 
            (\b{b}) not in the coefficient of variation in DNA content.
            These measurements are addressed fully in \ar{imaging:introduction}.}
  \label{fig:introduction:subpops}}
  \end{figure}



\subsubsection{Determining context-dependency}

Finally, how can we deal with the issue of context-dependency?
First, it is important to verify that the context-dependency truly exists.
As I discuss in \ar{pathways:introduction} and implied in this chapter,
context-dependency of biological phenomena is often inferred by
the fact that different labs produce different results when asking
the same questions. However, interpretation of cell signaling results
are incredibly complicated, which may
simply mean that the labs were not, in fact, asking the same questions.
Because cells
may be encoding information differently than we expect, we should
take care when comparing interpretations of results obtained by different
experimental methods, as each method will approximate $\vec{S}$ and $\vec{R}$
differently.


However, some (perhaps much) of context-dependency is undoubtedly
real, and can be absorbed into the simple model of cells as functions.
While we could allow the function to vary from context to context,
it is more useful to say that $f$ does not change but that
subsets of the inputs $\vec{S}$ and outputs $\vec{R}$ can vary.


To get around this parameter variation, we can first make sure that
all controllable conditions are kept the same and that all measurements
are the same from experiment to experiment. Thus, the experimentally-defined
subset of values in $\vec{S}$ and $\vec{R}$ do not change. Experiments can
then be repeated identically across multiple cell types, so that the only
varying parameters are those inherent to cell type differences.
Any consistent aspects of the relationship between experimental
$S$ and $R$ across diverse cell types can then be used to infer the
general properties of $f$. Indeed, this approach is common in cell biology,
as it is widely believed that any given cell line may have a myriad of
idiosyncratic properties.


When using such a multiple cell-type approach, one may find a case
where context-dependency is so dramatic that no
general properties of $f$ can be uncovered. The first aspect
of this problem to tackle would be to carefully ask if the cell types
are truly being treated ``identically.'' As suggested above,
experiments typically use an absolute set of conditions across cell types
(e.g. identical ligand or drug concentrations). But it may be the case that two cell
types simply vary in sensitivity to the conditions, such that one type
is effectively receiving a half-maximal dose while the other is saturated.
Because we do not know which property $R$ encodes the treatment condition,
it is also difficult to know if we are measuring an ``identical'' readout.
Perhaps some of the apparent context-dependency of signaling is due to
incorrectly interpreting what it means to treat different cell types
identically.


Instead of relying on constant treatment concentrations
derived from the literature and applying such conditions generally
across experiments, another approach would be to measure dose-response and time-response
curves for all cell types that are under experimentation. Conditions could then be
calibrated on a cell type-specific basis so that, for example, all cell
lines receive a half-maximal input concentration.


\subsubsection{The encoding problem is unsolved}

Part of the intention of this chapter was to make it clear that
cellular signaling is an incredibly difficult phenomenon to understand,
and that experimental designs are making many
assumptions that are either going unnoticed or are not
being made explicit. Some of these assumptions, if made explicit, might
dramatically affect how we interpret our experimental results.


There is no general solution to the encoding problem but,
as I have outlined here, steps can be taken to minimize its effects. Perhaps more
importantly, an awareness of the assumptions allows for them to be tested
in some cases or, at minimum, allows for results to be interpreted cautiously
in the light of those assumptions.

